% ==================================================================
% id:       func-008
% topic:    פונקציות
% subtopic: תחום הגדרה ותכונות של פונקציות
% solution: full
% ==================================================================
\textbf{מצאו את התמונה של הפונקציות הבאות:}
\begin{enumerate}\bMs{2}
\item $f(x) = \frac{2x}{x^2 + 1}$
\hint{כדי לבדוק אם ערך $y$ שייך לתמונה, פתרו את המשוואה $f(x)=y$ עבור $x$ ובדקו מתי יש לה פתרון ממשי.}
\begin{solbox}
הפונקציה מוגדרת לכל $x$ ממשי. מספר $y$ שייך לתמונה אם ורק אם למשוואה
$\frac{2x}{x^2+1}=y$
יש פתרון ממשי $x$. נכפול ב-$x^2+1>0$ ונקבל משוואה שקולה:
$$.yx^2-2x+y=0$$
אם $y=0$, הפתרון הוא $x=0$, ולכן $0$ בתמונה.
אם $y\neq0$, זו משוואה ריבועית ב-$x$, ויש לה פתרון ממשי אם ורק אם הדיסקרימיננטה אי-שלילית:
$$.4-4y^2\geq0\iff y^2\leq1\iff -1\leq y\leq1$$
לכן התמונה היא
$[-1, 1]$.
\end{solbox}
\item $f(x) = \sqrt{2 + x - x^2}$
\hint{השלימו לריבוע את הביטוי שתחת השורש, ובדקו בין אילו ערכים הוא נע בתחום ההגדרה.}
\begin{solbox}
נשלים לריבוע. קודם נוציא מינוס, ואז נשלים לריבוע את הביטוי שבסוגריים:
$$.2+x-x^2=-\lr{x^2-x-2}=-\lr{x^2-x+\frac14-\frac14-2}=-\lr{\lr{x-\frac12}^2-\frac94}=\frac94-\lr{x-\frac12}^2$$
הפונקציה מוגדרת כאשר הביטוי הזה אי-שלילי, כלומר כאשר
$\abs{x-\frac12}\leq\frac32$,
ולכן תחום ההגדרה הוא $[-1,2]$.

לכל $x$ בתחום מתקיים
$0\leq \frac94-\lr{x-\frac12}^2\leq\frac94$,
ולכן
$$.0\leq f(x)\leq\frac32$$
נראה שכל ערך $y$ בקטע $[0,\frac32]$ מתקבל. נבחר
$x=\frac12+\sqrt{\frac94-y^2}$.
זהו מספר בקטע $[\frac12,2]$, ולכן הוא בתחום ההגדרה, ומתקיים
$$.f(x)=\sqrt{\frac94-\lr{\frac94-y^2}}=\sqrt{y^2}=y$$
לכן התמונה היא
$[0, \frac{3}{2}]$.
\end{solbox}
\eMs
\end{enumerate}
